%
\begin{isabellebody}%
\setisabellecontext{Week{\isacharunderscore}{\kern0pt}{\isadigit{2}}{\isacharunderscore}{\kern0pt}sol}%
%
\isadelimtheory
\isanewline
%
\endisadelimtheory
%
\isatagtheory
\isacommand{theory}\isamarkupfalse%
\ Week{\isacharunderscore}{\kern0pt}{\isadigit{2}}{\isacharunderscore}{\kern0pt}sol\isanewline
\isakeyword{imports}\ Main\isanewline
\isakeyword{begin}%
\endisatagtheory
{\isafoldtheory}%
%
\isadelimtheory
%
\endisadelimtheory
%
\begin{isamarkuptext}%
\vspace{15ex}\ExerciseSheet{7}{}%
\end{isamarkuptext}\isamarkuptrue%
%
\begin{isamarkuptext}%
\Exercise{Fold function}
  The fold function is a very generic function, that can be used to express 
  multiple other interesting functions over lists.%
\end{isamarkuptext}\isamarkuptrue%
%
\begin{isamarkuptext}%
Have a look at Isabelle/HOL's standard function \isa{fold}.%
\end{isamarkuptext}\isamarkuptrue%
\isacommand{thm}\isamarkupfalse%
\ fold{\isachardot}{\kern0pt}simps%
\begin{isamarkuptext}%
Write a function to compute the sum of the elements of a list. 
  Define two versions, one direct recursive definition, and one using fold.
  Show that both are equal.

  Hint: use automation!%
\end{isamarkuptext}\isamarkuptrue%
\isacommand{fun}\isamarkupfalse%
\ list{\isacharunderscore}{\kern0pt}sum\ {\isacharcolon}{\kern0pt}{\isacharcolon}{\kern0pt}\ {\isachardoublequoteopen}nat\ list\ {\isasymRightarrow}\ nat{\isachardoublequoteclose}\ \isanewline
\ \ \ \ \isanewline
\ \ \isakeyword{where}\isanewline
\ \ \ \ {\isachardoublequoteopen}list{\isacharunderscore}{\kern0pt}sum\ {\isacharbrackleft}{\kern0pt}{\isacharbrackright}{\kern0pt}\ {\isacharequal}{\kern0pt}\ {\isadigit{0}}{\isachardoublequoteclose}\isanewline
\ \ {\isacharbar}{\kern0pt}\ {\isachardoublequoteopen}list{\isacharunderscore}{\kern0pt}sum\ {\isacharparenleft}{\kern0pt}x{\isacharhash}{\kern0pt}xs{\isacharparenright}{\kern0pt}\ {\isacharequal}{\kern0pt}\ x\ {\isacharplus}{\kern0pt}\ list{\isacharunderscore}{\kern0pt}sum\ xs{\isachardoublequoteclose}\ \ \isanewline
\ \ \ \ \ \ \isanewline
\isanewline
\isacommand{definition}\isamarkupfalse%
\ list{\isacharunderscore}{\kern0pt}sum{\isacharprime}{\kern0pt}\ {\isacharcolon}{\kern0pt}{\isacharcolon}{\kern0pt}\ {\isachardoublequoteopen}nat\ list\ {\isasymRightarrow}\ nat{\isachardoublequoteclose}\isanewline
\ \ \ \ \isanewline
\ \ \isakeyword{where}\ {\isachardoublequoteopen}list{\isacharunderscore}{\kern0pt}sum{\isacharprime}{\kern0pt}\ xs\ {\isasymequiv}\ fold\ {\isacharparenleft}{\kern0pt}{\isacharplus}{\kern0pt}{\isacharparenright}{\kern0pt}\ xs\ {\isadigit{0}}{\isachardoublequoteclose}\isanewline
\ \ \ \ \ \ \ \ \isanewline
\isanewline
\ \ \ \ \isanewline
\isacommand{lemma}\isamarkupfalse%
\ aux{\isacharcolon}{\kern0pt}\ {\isachardoublequoteopen}fold\ {\isacharparenleft}{\kern0pt}{\isacharplus}{\kern0pt}{\isacharparenright}{\kern0pt}\ xs\ a\ {\isacharequal}{\kern0pt}\ list{\isacharunderscore}{\kern0pt}sum\ xs\ {\isacharplus}{\kern0pt}\ a{\isachardoublequoteclose}\ \ \isanewline
%
\isadelimproof
\ \ %
\endisadelimproof
%
\isatagproof
\isacommand{by}\isamarkupfalse%
\ {\isacharparenleft}{\kern0pt}induction\ xs\ arbitrary{\isacharcolon}{\kern0pt}\ a{\isacharparenright}{\kern0pt}\ auto%
\endisatagproof
{\isafoldproof}%
%
\isadelimproof
\isanewline
%
\endisadelimproof
\ \ \ \ \ \ \ \ \isanewline
\isanewline
\isacommand{lemma}\isamarkupfalse%
\ {\isachardoublequoteopen}list{\isacharunderscore}{\kern0pt}sum\ xs\ {\isacharequal}{\kern0pt}\ list{\isacharunderscore}{\kern0pt}sum{\isacharprime}{\kern0pt}\ xs{\isachardoublequoteclose}\isanewline
%
\isadelimproof
\ \ \ \ \ \ \isanewline
\ \ %
\endisadelimproof
%
\isatagproof
\isacommand{unfolding}\isamarkupfalse%
\ list{\isacharunderscore}{\kern0pt}sum{\isacharprime}{\kern0pt}{\isacharunderscore}{\kern0pt}def\ \isacommand{by}\isamarkupfalse%
\ {\isacharparenleft}{\kern0pt}simp\ add{\isacharcolon}{\kern0pt}\ aux{\isacharparenright}{\kern0pt}%
\endisatagproof
{\isafoldproof}%
%
\isadelimproof
%
\endisadelimproof
%
\begin{isamarkuptext}%
\Exercise{Tail recursive \isa{reverse}}
In Isabelle, there is the reverse function \isa{rev}, which, given a list, computes another list
with the same elements but in reverse order.
Take a loot into its definition:%
\end{isamarkuptext}\isamarkuptrue%
\isacommand{thm}\isamarkupfalse%
\ rev{\isachardot}{\kern0pt}simps%
\begin{isamarkuptext}%
Recall tail recursion: a function is tail recursive if in all recursive calls it does not call 
  anything after itself\footnote{See \url{https://en.wikipedia.org/wiki/Tail_call.} for a quick review.}.
  For instance, the implementaion of \isa{rev} in Isabelle is not tail recursive, because it calls
  \isa{{\isacharparenleft}{\kern0pt}{\isacharat}{\kern0pt}{\isacharparenright}{\kern0pt}} after it calls \isa{rev} in the second equation.
  
  Write an alternative implementation of the reverse function that is tail recursive.

  Hint: recall that you can most of the times derive a tail recursive function using an accumulator
        argument.%
\end{isamarkuptext}\isamarkuptrue%
\isacommand{fun}\isamarkupfalse%
\ rev{\isacharprime}{\kern0pt}{\isacharcolon}{\kern0pt}{\isacharcolon}{\kern0pt}{\isachardoublequoteopen}{\isacharprime}{\kern0pt}a\ list\ {\isasymRightarrow}\ {\isacharprime}{\kern0pt}a\ list\ {\isasymRightarrow}\ {\isacharprime}{\kern0pt}a\ list{\isachardoublequoteclose}\ \isakeyword{where}\isanewline
\ \ {\isachardoublequoteopen}rev{\isacharprime}{\kern0pt}\ {\isacharbrackleft}{\kern0pt}{\isacharbrackright}{\kern0pt}\ acc\ {\isacharequal}{\kern0pt}\ acc{\isachardoublequoteclose}\isanewline
{\isacharbar}{\kern0pt}\ {\isachardoublequoteopen}rev{\isacharprime}{\kern0pt}\ {\isacharparenleft}{\kern0pt}x\ {\isacharhash}{\kern0pt}xs{\isacharparenright}{\kern0pt}\ acc\ {\isacharequal}{\kern0pt}\ rev{\isacharprime}{\kern0pt}\ xs\ {\isacharparenleft}{\kern0pt}x\ {\isacharhash}{\kern0pt}\ acc{\isacharparenright}{\kern0pt}{\isachardoublequoteclose}\isanewline
\isanewline
\isanewline
\isanewline
\isanewline
\isacommand{fun}\isamarkupfalse%
\ rev{\isacharunderscore}{\kern0pt}tr{\isacharcolon}{\kern0pt}{\isacharcolon}{\kern0pt}{\isachardoublequoteopen}{\isacharprime}{\kern0pt}a\ list\ {\isasymRightarrow}\ {\isacharprime}{\kern0pt}a\ list{\isachardoublequoteclose}\ \isakeyword{where}\isanewline
\isanewline
\ \ {\isachardoublequoteopen}rev{\isacharunderscore}{\kern0pt}tr\ xs\ {\isacharequal}{\kern0pt}\ rev{\isacharprime}{\kern0pt}\ xs\ {\isacharbrackleft}{\kern0pt}{\isacharbrackright}{\kern0pt}{\isachardoublequoteclose}%
\begin{isamarkuptext}%
Show that the two implementations are equivalent.%
\end{isamarkuptext}\isamarkuptrue%
\isacommand{lemma}\isamarkupfalse%
\ rev{\isacharprime}{\kern0pt}{\isacharunderscore}{\kern0pt}append{\isacharcolon}{\kern0pt}\ {\isachardoublequoteopen}rev{\isacharprime}{\kern0pt}\ xs\ ys\ {\isacharequal}{\kern0pt}\ {\isacharparenleft}{\kern0pt}rev{\isacharprime}{\kern0pt}\ xs\ {\isacharbrackleft}{\kern0pt}{\isacharbrackright}{\kern0pt}{\isacharparenright}{\kern0pt}\ {\isacharat}{\kern0pt}\ ys{\isachardoublequoteclose}\isanewline
%
\isadelimproof
%
\endisadelimproof
%
\isatagproof
\isacommand{proof}\isamarkupfalse%
\ {\isacharparenleft}{\kern0pt}induction\ xs\ arbitrary{\isacharcolon}{\kern0pt}\ ys{\isacharparenright}{\kern0pt}\isanewline
\ \ \isacommand{case}\isamarkupfalse%
\ Nil\isanewline
\ \ \isacommand{then}\isamarkupfalse%
\ \isacommand{show}\isamarkupfalse%
\ {\isacharquery}{\kern0pt}case\ \isanewline
\ \ \ \ \isacommand{by}\isamarkupfalse%
\ auto\isanewline
\isacommand{next}\isamarkupfalse%
\isanewline
\ \ \isacommand{case}\isamarkupfalse%
\ {\isacharparenleft}{\kern0pt}Cons\ a\ xs{\isacharparenright}{\kern0pt}\isanewline
\ \ \isacommand{have}\isamarkupfalse%
\ {\isadigit{1}}{\isacharcolon}{\kern0pt}\ {\isachardoublequoteopen}rev{\isacharprime}{\kern0pt}\ xs\ {\isacharparenleft}{\kern0pt}a\ {\isacharhash}{\kern0pt}\ ys{\isacharparenright}{\kern0pt}\ {\isacharequal}{\kern0pt}\ {\isacharparenleft}{\kern0pt}rev{\isacharprime}{\kern0pt}\ xs\ {\isacharbrackleft}{\kern0pt}{\isacharbrackright}{\kern0pt}{\isacharparenright}{\kern0pt}\ {\isacharat}{\kern0pt}\ {\isacharparenleft}{\kern0pt}a\ {\isacharhash}{\kern0pt}\ ys{\isacharparenright}{\kern0pt}{\isachardoublequoteclose}\isanewline
\ \ \ \ \isacommand{apply}\isamarkupfalse%
{\isacharparenleft}{\kern0pt}subst\ Cons{\isacharparenright}{\kern0pt}\isanewline
\ \ \ \ \isacommand{{\isachardot}{\kern0pt}{\isachardot}{\kern0pt}}\isamarkupfalse%
\isanewline
\ \ \isacommand{have}\isamarkupfalse%
\ {\isadigit{2}}{\isacharcolon}{\kern0pt}\ {\isachardoublequoteopen}rev{\isacharprime}{\kern0pt}\ xs\ {\isacharbrackleft}{\kern0pt}a{\isacharbrackright}{\kern0pt}\ {\isacharequal}{\kern0pt}\ {\isacharparenleft}{\kern0pt}rev{\isacharprime}{\kern0pt}\ xs\ {\isacharbrackleft}{\kern0pt}{\isacharbrackright}{\kern0pt}{\isacharparenright}{\kern0pt}\ {\isacharat}{\kern0pt}\ {\isacharbrackleft}{\kern0pt}a{\isacharbrackright}{\kern0pt}{\isachardoublequoteclose}\isanewline
\ \ \ \ \isacommand{apply}\isamarkupfalse%
{\isacharparenleft}{\kern0pt}subst\ Cons{\isacharparenright}{\kern0pt}\isanewline
\ \ \ \ \isacommand{{\isachardot}{\kern0pt}{\isachardot}{\kern0pt}}\isamarkupfalse%
\isanewline
\ \ \isacommand{show}\isamarkupfalse%
\ {\isacharquery}{\kern0pt}case\isanewline
\ \ \ \ \isacommand{using}\isamarkupfalse%
\ {\isadigit{1}}\ {\isadigit{2}}\isanewline
\ \ \ \isacommand{by}\isamarkupfalse%
\ simp\isanewline
\isacommand{qed}\isamarkupfalse%
%
\endisatagproof
{\isafoldproof}%
%
\isadelimproof
\isanewline
%
\endisadelimproof
\isanewline
\ \isanewline
\isacommand{lemma}\isamarkupfalse%
\ {\isachardoublequoteopen}rev{\isacharunderscore}{\kern0pt}tr\ xs\ {\isacharequal}{\kern0pt}\ rev\ xs{\isachardoublequoteclose}\isanewline
%
\isadelimproof
\isanewline
%
\endisadelimproof
%
\isatagproof
\isacommand{proof}\isamarkupfalse%
{\isacharparenleft}{\kern0pt}induction\ xs{\isacharparenright}{\kern0pt}\isanewline
\ \ \isacommand{case}\isamarkupfalse%
\ Nil\isanewline
\ \ \isacommand{then}\isamarkupfalse%
\ \isacommand{show}\isamarkupfalse%
\ {\isacharquery}{\kern0pt}case\isanewline
\ \ \ \ \isacommand{by}\isamarkupfalse%
\ auto\isanewline
\isacommand{next}\isamarkupfalse%
\isanewline
\ \ \isacommand{case}\isamarkupfalse%
\ {\isacharparenleft}{\kern0pt}Cons\ a\ xs{\isacharparenright}{\kern0pt}\isanewline
\ \ \isacommand{then}\isamarkupfalse%
\ \isacommand{show}\isamarkupfalse%
\ {\isacharquery}{\kern0pt}case\isanewline
\ \ \ \ \isacommand{apply}\isamarkupfalse%
\ simp\isanewline
\ \ \ \ \isacommand{apply}\isamarkupfalse%
{\isacharparenleft}{\kern0pt}subst\ rev{\isacharprime}{\kern0pt}{\isacharunderscore}{\kern0pt}append{\isacharparenright}{\kern0pt}\isanewline
\ \ \ \ \isacommand{by}\isamarkupfalse%
\ auto\isanewline
\isacommand{qed}\isamarkupfalse%
%
\endisatagproof
{\isafoldproof}%
%
\isadelimproof
\isanewline
%
\endisadelimproof
\isanewline
\isanewline
%
\isadelimtheory
\isanewline
%
\endisadelimtheory
%
\isatagtheory
\isacommand{end}\isamarkupfalse%
%
\endisatagtheory
{\isafoldtheory}%
%
\isadelimtheory
%
\endisadelimtheory
%
\end{isabellebody}%
\endinput
%:%file=Week_2_sol.tex%:%
%:%6=1%:%
%:%11=2%:%
%:%12=2%:%
%:%13=3%:%
%:%14=4%:%
%:%23=6%:%
%:%27=8%:%
%:%28=9%:%
%:%29=10%:%
%:%33=13%:%
%:%35=14%:%
%:%36=14%:%
%:%38=17%:%
%:%39=18%:%
%:%40=19%:%
%:%41=20%:%
%:%42=21%:%
%:%44=26%:%
%:%45=26%:%
%:%46=27%:%
%:%47=28%:%
%:%48=29%:%
%:%49=30%:%
%:%50=31%:%
%:%51=32%:%
%:%52=33%:%
%:%53=33%:%
%:%54=34%:%
%:%55=35%:%
%:%56=36%:%
%:%57=37%:%
%:%58=38%:%
%:%59=39%:%
%:%60=39%:%
%:%63=40%:%
%:%67=40%:%
%:%68=40%:%
%:%73=40%:%
%:%76=41%:%
%:%77=42%:%
%:%78=43%:%
%:%79=43%:%
%:%82=44%:%
%:%83=45%:%
%:%87=45%:%
%:%88=45%:%
%:%89=45%:%
%:%98=48%:%
%:%99=49%:%
%:%100=50%:%
%:%101=51%:%
%:%103=54%:%
%:%104=54%:%
%:%106=57%:%
%:%107=58%:%
%:%108=59%:%
%:%109=60%:%
%:%110=61%:%
%:%111=62%:%
%:%112=63%:%
%:%113=64%:%
%:%114=65%:%
%:%116=70%:%
%:%117=70%:%
%:%118=71%:%
%:%119=72%:%
%:%120=73%:%
%:%121=74%:%
%:%122=75%:%
%:%123=76%:%
%:%124=77%:%
%:%125=77%:%
%:%126=78%:%
%:%127=79%:%
%:%129=82%:%
%:%131=87%:%
%:%132=87%:%
%:%139=88%:%
%:%140=88%:%
%:%141=89%:%
%:%142=89%:%
%:%143=90%:%
%:%144=90%:%
%:%145=90%:%
%:%146=91%:%
%:%147=91%:%
%:%148=92%:%
%:%149=92%:%
%:%150=93%:%
%:%151=93%:%
%:%152=94%:%
%:%153=94%:%
%:%154=95%:%
%:%155=95%:%
%:%156=96%:%
%:%157=96%:%
%:%158=97%:%
%:%159=97%:%
%:%160=98%:%
%:%161=98%:%
%:%162=99%:%
%:%163=99%:%
%:%164=100%:%
%:%165=100%:%
%:%166=101%:%
%:%167=101%:%
%:%168=102%:%
%:%169=102%:%
%:%170=103%:%
%:%176=103%:%
%:%179=104%:%
%:%180=105%:%
%:%181=106%:%
%:%182=106%:%
%:%185=107%:%
%:%190=108%:%
%:%191=108%:%
%:%192=109%:%
%:%193=109%:%
%:%194=110%:%
%:%195=110%:%
%:%196=110%:%
%:%197=111%:%
%:%198=111%:%
%:%199=112%:%
%:%200=112%:%
%:%201=113%:%
%:%202=113%:%
%:%203=114%:%
%:%204=114%:%
%:%205=114%:%
%:%206=115%:%
%:%207=115%:%
%:%208=116%:%
%:%209=116%:%
%:%210=117%:%
%:%211=117%:%
%:%212=118%:%
%:%218=118%:%
%:%221=119%:%
%:%222=120%:%
%:%225=121%:%
%:%230=122%:%
